\section{Introduction}\label{introduction}

Does a language model that behaves morally understand morality, or has it
only learned to comply? Earlier papers in this series studied what models
\emph{represent}: moral content is linearly decodable early in pre-training and
across nearly all layers \citep{reblitzrichardson2026fragility}, is encoded
uniformly across mixture-of-experts experts rather than specialized
\citep{reblitzrichardson2026dilution}, and organizes into structured,
five-dimensional framework geometry whose directions are causal for moral
judgments \citep{reblitzrichardson2026geometry, reblitzrichardson2026causal}.
All of that was measured on base, pre-trained models. The moral behavior we
actually deploy and align, refusing harmful requests and giving
morally-aware answers, is produced by \emph{post-training}: supervised fine-tuning
(SFT), preference optimization (DPO), and reinforcement learning with
verifiable rewards (RLVR). This paper asks what post-training does to moral
representation. Does alignment \emph{teach} morality, or does it \emph{wire} a
pre-existing moral representation to behavior?

We separate two things a model can have. \textbf{Comprehension} is the moral content
a model internally represents: which foundation a scenario engages, how
foundations relate, whether the representation is decodable and stable.
\textbf{Compliance} is whether the model's behavior conforms: whether it refuses a
harmful request, whether its stated judgment matches the morally expected one.
A model can have either without the other. Our hypothesis is that alignment is a
\textbf{coupling change, not a teaching change}: post-training does not create moral
understanding, it connects understanding that pre-training already built to
behavioral output. If so, failures of alignment, jailbreaks and ablations, are
failures of \emph{coupling}, not of \emph{understanding}.

We test this on the full OLMo-3 7B pipeline, which publishes every stage:
the base model and 13 stage-3 pre-training checkpoints, the SFT and DPO
snapshots, and the final RLVR-tuned Instruct model with 8 intermediate
checkpoints, 25 model states in all, plus a refusal-ablated instruct.
Across them we measure comprehension (foundation probe-direction transfer and
geometry), compliance (behavioral moral judgment and refusal), the coupling
between them, the persona direction and its angle to the moral subspace, and
the geometry of the refusal direction itself. The results form a chain, in
which each finding raises the question the next one answers.

\textbf{Finding 1: Moral comprehension is a pre-training property.} Across all 25
states the six foundations are linearly decodable at 100\%, span five effective
dimensions, and base-trained directions transfer as near-perfect classifiers
(AUC \(\approx 1.0\)). The moral subspace \emph{crystallizes} during pre-training:
the cosine between a checkpoint's foundation directions and the final base
model's rises monotonically from \(0.87\) to \(0.999\) over the stage-3 anneal.
Comprehension is in place before any alignment.

If comprehension is already present at the base, the natural question is what
post-training changes. \textbf{Finding 2: post-training reorients moral
representation, it does not re-teach it, and SFT does so once.} SFT applies a
single \({\sim}40^{\circ}\) rotation to the moral subspace (direction cosine to base
\(0.999 \rightarrow 0.757\)); DPO and all eight RLVR checkpoints leave it flat.
Effective dimensionality stays at five, mean pairwise foundation cosine barely
moves (\(0.262 \rightarrow 0.250\)), and the framework clustering is preserved.
Alignment rotates the moral representation slightly and then leaves it alone.

If post-training preserves comprehension rather than building it, then whatever
alignment adds to behavior must be coupled to that fixed representation only
loosely or not at all. \textbf{Finding 3: comprehension and compliance are only
weakly coupled, and that weak coupling is itself the result.} Per-scenario
agreement between the model's internally dominant foundation and its behavioral
correctness rises across post-training (\(0.375 \rightarrow 0.500\); \(\phi\)
\(-0.19 \rightarrow +0.05\)) but remains weak: even at the final Instruct model,
\(P(\text{comply} \mid \text{comprehend}) \approx P(\text{comply} \mid
\neg\text{comprehend}) \approx 0.75\). A near-zero coupling is not a null result;
it states that in the fully aligned model, moral representation and behavioral
compliance are barely linked. This predicts a specific, testable consequence:
if compliance were routed \emph{through} moral representations, removing compliance
would damage comprehension; weak coupling predicts instead that compliance is a
\emph{separate} mechanism that can be removed while comprehension survives.

\textbf{Finding 4: the refusal mechanism is geometrically separate from morality, and
removing it confirms the dissociation.} The refusal direction recovered by
Heretic-style difference-of-means \citep{arditi2024refusal} is nearly orthogonal
to the six-foundation moral subspace (projection fraction \(0.10\), mean
\(|\cos| = 0.06\)). Ablating it leaves comprehension untouched (effective
dimensionality 5, direction transfer and probe accuracy identical to the
un-ablated instruct) and moral judgment intact (behavioral accuracy
\(0.75 \rightarrow 0.73\)), while refusal collapses (refusal rate
\(0.25 \rightarrow 0.00\)): the ablated model now answers requests it previously
refused. This is the high-comprehension, low-compliance cell of the
dissociation, realized in a real model.

\textbf{The separation is a vulnerability, and it cannot be repaired after
pre-training.} If compliance is a separable mechanism, the obvious defense is to
couple it back to comprehension, so that the refusal-direction ablation of
Finding 4 can no longer remove one without the other. We first measure how much
the model relies on its moral subspace for generation, a reliance that is
present at the base and grows roughly sixfold across alignment, then attempt to
deepen it with ablation-resistance training: an auxiliary fine-tuning loss that
rewards the model when projecting the moral subspace out of its activations
degrades its output. The intervention builds strong dependence on the moral
subspace but does not couple compliance to it. That dependence is orthogonal to
the refusal direction an attacker removes, and refusal ablation damages the
trained model no more than the control. The failure is developmental: the
refusal feature forms during post-training out of pre-training proto-features
that are already orthogonal to the moral subspace, so a post-training loss acts
too late to relocate it. Making alignment ablation-resistant is, on this
evidence, a pre-training problem.

Together these findings populate a comprehension\(\times\)compliance matrix in
which the low-comprehension row is empty, every OLMo-3 state has full moral
decodability, and all variation lives in the high-comprehension row, where
compliance moves independently of understanding. The contribution is twofold: a
method for measuring comprehension, compliance, and their coupling across an
alignment pipeline, with the result that alignment is a wiring step that
preserves moral comprehension built in pre-training and attaches a separable
compliance mechanism, so that compliance can be removed without touching what the
model understands; and the finding that this separation cannot be undone after
the fact, which locates the problem of ablation-resistant alignment in
pre-training rather than post-training.
